\documentclass[12pt]{article}
\usepackage[T1]{fontenc}
\usepackage[utf8]{inputenc}
\usepackage{lmodern}
\usepackage{textcomp}
\usepackage{enumitem}
\usepackage{geometry}
\geometry{margin=1in}
\begin{document}
\begin{center}
Answer Key - Physics - Graduate Level Assessment 18 \
\small (Answers are ordered exactly as in the question paper.)
\end{center}

\section*{Multiple Choice}
\begin{enumerate}[label=\arabic*.]
\item \textbf{Answer:} D
\\ \textit{Options:} A) Neptune; B) Uranus; C) Saturn; D) all of the above
\\ \textit{Note:} All of the jovian planets-Jupiter, Saturn, Uranus, and Neptune-possess ring systems, though their visibility and composition vary significantly, confirming that each has rings.
\item \textbf{Answer:} A
\\ \textit{Options:} A) The Moon and Earth have nearly identical isotopic compositions.; B) The Moon's orbit is not in the same plane as Earth's equator.; C) The Moon's crust is thicker on the far side than the near side.; D) The Moon has a much smaller iron core than the Earth even considering its size.; E) The Moon was once entirely molten.
\\ \textit{Note:} The Giant Impact Hypothesis suggests that the Moon was formed from debris resulting from a collision between the early Earth and a Mars-sized body, and the nearly identical isotopic compositions of the Earth and Moon indicate a shared origin rather than an external source, which contradicts the idea of it being evidence for the hypothesis.
\item \textbf{Answer:} B
\\ \textit{Options:} A) 750.0; B) 924.0; C) 1300.0; D) 800.0; E) 870.0
\\ \textit{Note:} The correct answer of 924.0 volts is derived from applying the relativistic energy-momentum relationship, where the kinetic energy gained by the electron when accelerated through a potential difference equates to its relativistic kinetic energy, leading to the calculation of the necessary voltage to achieve the specified speed of 1.8 x $10^7$ m/s.
\item \textbf{Answer:} A
\\ \textit{Options:} A) frequency does not change, the wavelength increases, and the speed increases; B) frequency decreases, the wavelength increases, and the speed does not change; C) frequency increases, the wavelength increases, and the speed decreases; D) frequency increases, the wavelength does not change, and the speed increases; E) frequency increases, the wavelength decreases, and the speed does not change
\\ \textit{Note:} When a wave transitions from shallow water to deep water, the frequency remains constant while the wavelength and speed increase due to the deeper water allowing for faster wave propagation, as described by the wave speed equation v = fλ, where v is speed, f is frequency, and λ is wavelength.
\item \textbf{Answer:} D
\\ \textit{Options:} A) 1000 Mev; B) 650 MeV; C) 800 MeV; D) 931 Mev; E) 500 Mev
\\ \textit{Note:} The correct answer is 931 MeV because 1 atomic mass unit (amu) is equivalent to approximately 931 million electron volts, which is derived from the mass-energy equivalence principle articulated by Einstein's equation E=mc², where the mass of 1 amu can be converted into energy.
\end{enumerate}

\section*{True / False}
\begin{enumerate}[label=\arabic*.]
\setcounter{enumi}{5}
\item \textbf{Answer:} True
\\ \textit{Options:} A) True; B) False
\\ \textit{Note:} The energy of a photon is inversely proportional to its wavelength, and using the formula E = hc/λ, where h is Planck's constant and c is the speed of light, the energy corresponding to a wavelength of 400 nm is calculated to be approximately 3.1 eV.
\item \textbf{Answer:} False
\\ \textit{Options:} A) True; B) False
\\ \textit{Note:} The value of g in Paris, which is approximately 9.81 m/s² or 32.2 ft/s², cannot be approximated as 33.000 ft/s² because this figure exceeds the actual acceleration due to gravity experienced in that location.
\item \textbf{Answer:} True
\\ \textit{Options:} A) True; B) False
\\ \textit{Note:} The statement is true because the solar cycle, which reflects the periodic increase and decrease in solar activity, including sunspot numbers, has an average duration of approximately 11 years due to the complex interactions of the Sun's magnetic field.
\item \textbf{Answer:} False
\\ \textit{Options:} A) True; B) False
\\ \textit{Note:} The statement is false because the deposition of silver can be calculated using Faraday's laws of electrolysis, which show that the amount of substance deposited is proportional to the current and time, and for silver, 0.50 A for 20 minutes would yield less than 1.000 grams.
\item \textbf{Answer:} False
\\ \textit{Options:} A) True; B) False
\\ \textit{Note:} The time required for a satellite to complete one revolution, known as the orbital period, is determined by the radius of its orbit and the mass of the planet only affects the gravitational force, not the proportionality to the orbital period itself, which follows Kepler's third law of planetary motion.
\end{enumerate}

\section*{Long Form Response}
\begin{enumerate}[label=\arabic*.]
\setcounter{enumi}{10}
\item \textbf{Answer:} The U.S.S.R.'s lander missions to Venus, particularly the Venera program, were pivotal in advancing planetary exploration by providing direct data from the planet's surface and atmosphere. These missions, which successfully landed several probes on Venus, delivered crucial information about the planet's high temperatures, pressure, and chemical composition, revealing a hostile environment dominated by carbon dioxide and sulfuric acid clouds. The geological data collected helped scientists understand the planet's volcanic activity and surface features, suggesting a younger geological age than previously thought. Furthermore, the Venera landers offered insights into the atmospheric dynamics of Venus, contributing to broader theories regarding planetary atmospheres and climate. Overall, these missions not only expanded our knowledge of Venus but also established methodologies for future planetary exploration.
\\ \textit{Note:} correct because it emphasizes the Venera program's role in gathering direct, vital data on Venus's extreme environmental conditions and geological characteristics, significantly enhancing our understanding of the planet within the broader context of planetary exploration.
\item \textbf{Answer:} The work done by the 5 N horizontal force over 1 minute, which is 60 seconds, can be calculated using the formula for work: \( W = F \times d \), where \( d \) is the distance traveled. Since the box moves at a constant velocity of 2 m/s, the distance covered in one minute is \( d = v \times t = 2 \, \text{m/s} \times 60 \, \text{s} = 120 \, \text{m} \). Substituting the values gives \( W = 5 \, \text{N} \times 120 \, \text{m} = 600 \, \text{J} \). This scenario highlights the role of friction, as the force of 5 N is equal to the frictional force opposing the motion, allowing the box to maintain a constant velocity rather than accelerating. Thus, the work done by the applied force is entirely used to overcome friction, illustrating the balance of forces in a system at equilibrium.
\\ \textit{Note:} correct because it accurately applies the work formula by calculating the distance traveled at a constant velocity under the influence of a horizontal force, while also implicitly recognizing that the force must counteract friction to maintain that constant velocity.
\end{enumerate}

\vfill
\noindent\textit{End of Answer Key}
\end{document}